% Generated draft: requires mathematical and source review.
\documentclass[11pt]{article}
\usepackage[margin=25mm]{geometry}
\usepackage{fontspec}
\setmainfont{Latin Modern Roman}
\usepackage{amsmath,amssymb,mathtools,unicode-math}
\setmathfont{Latin Modern Math}
\usepackage{enumitem,needspace,xurl,hyperref}
\setlength{\parindent}{0pt}
\setlength{\parskip}{5pt}
\setlength{\emergencystretch}{3em}
\newcommand{\N}{\mathbb{N}}
\newcommand{\Z}{\mathbb{Z}}
\newcommand{\Q}{\mathbb{Q}}
\newcommand{\R}{\mathbb{R}}
\newcommand{\C}{\mathbb{C}}
\newcommand{\F}{\mathbb{F}}
\newcommand{\A}{\mathbb{A}}
\newcommand{\PP}{\mathbb{P}}
\newcommand{\E}{\mathbb{E}}
\newcommand{\eps}{\varepsilon}
\newcommand{\dd}{\,\mathrm{d}}
\begin{document}
\section*{UnsolvedMath ID: 700208}
\textbf{Source problem code:} catalogue entry
\subsection{Definitions and mathematical statement}
Let \(G\) be a finite group, let \(p\) be a prime, and let \(P\) be a Sylow \(p\)-subgroup of \(G\): if \(|G|=p^a m\), where \(a\geq 0\) and \(p\nmid m\), then \(|P|=p^a\).

Write \(\ell_p(G)\) provisionally for the notion of \(p\)-length of an arbitrary finite group attributed in the source to L. A. Shemetkov, \emph{Math. USSR Sbornik}, \textbf{1} (1968), 83--92. Write \(c_p(P),d_p(P),e_p(P)\) for the three invariants of the Sylow subgroup mentioned in the source. These symbols are placeholders: the supplied record gives neither the definition of this generalized \(p\)-length nor the definitions of these invariants.

\textbf{Research problem (L. A. Shemetkov, 1969).} Investigate the relations between \(\ell_p(G)\) and \(c_p(P),d_p(P),e_p(P)\) for arbitrary finite groups \(G\).

The source poses an investigation, without specifying an inequality, equivalence, or other proposition to establish. It does not restrict \(G\) to be \(p\)-soluble, so a definition of \(p\)-length confined to that class would not specify the stated problem. The intended definitions and the conventions when \(P\) is trivial must be recovered before a fully precise formulation can be given; the source does not explicitly restrict to primes dividing \(|G|\).
\subsection{Short English statement}
Shemetkov asks how his notion of \(p\)-length for an arbitrary finite group relates to three invariants \(c_p,d_p,e_p\) of its Sylow \(p\)-subgroup. The supplied record does not define these quantities or specify a particular relation to prove, so this remains an underdetermined research question.
\subsection{Sources}
\begin{itemize}
\item Local UnsolvedMath registry, supplied complete source record: permanent ID \(700208\), exact source problem code \texttt{catalogue entry}, title \emph{Kourovka Notebook Problem 3.60}; upstream revision \texttt{d371ee917ac9647fbb06cd6fc529deec6f3dda82}. The registry labels the problem \texttt{open}; this is supplied status metadata, not an independently verified mathematical assessment.
\item \emph{The Kourovka Notebook (Unsolved Problems in Group Theory)}, Problem 3.60, attributed to L. A. Shemetkov, 3rd issue (1969). The supplied record identifies the 21st issue edition (2026), \texttt{arXiv:1401.0300v48}; its collection description names E. I. Khukhro and V. D. Mazurov as editors.
\item L. A. Shemetkov, \emph{Math. USSR Sbornik}, \textbf{1} (1968), 83--92, cited by the supplied statement for the definition of \(p\)-length of an arbitrary finite group. The record supplies no article title or definition. Bibliographic details are reproduced as supplied and have not been independently verified.
\end{itemize}
\end{document}
